% %%%%%%%%%%%%%%%%%%%%%%%%%%%%%%%%%%%%%%%%%%%%%%%%%%%%%%%%%%%%%%%%%%%%%%%%%%%%
% Thesis Defence
% Section: Conclusion
% %%%%%%%%%%%%%%%%%%%%%%%%%%%%%%%%%%%%%%%%%%%%%%%%%%%%%%%%%%%%%%%%%%%%%%%%%%%%

\section[Conclusion]{Conclusion}

\begin{frame}
    \frametitle{Conclusion}

    \textbf{\textcolor{Mulberry}{Objectif 1 :} Réduire l'impact de l’irrégularité des réductions en analysant le comportement des caches de données}

    \smallskip
    
    \arrow[\bfseries\color{Red}][2]{SpDCache}

    \point[][6][Red]{Cache de réduction}
    
    \point[\bfseries][6][Red]{Régions mémoires} du cache pour les régions \textbf{denses} et \textbf{creuses} des matrices

    \point[][6][Red]{Architecture \textbf{peu coûteuse} et \textbf{flexible}}

    \point[\bfseries][6][Red]{Implémentation matérielle} (RTL)

    \subpt[][10][Green]{\parbox[t]{3cm}{Traffic mémoire :}} \textbf{\textcolor{Green}{3.9$\times$}}

    \subpt[][10][Green]{\parbox[t]{3cm}{Accélération :}} \textbf{\textcolor{Green}{10\%}}
    
    \bigskip
    
    \textbf{\textcolor{Mulberry}{Objectif 2 :} Comprendre le rôle des structures des matrices dans les performances}

    \smallskip

    \arrow[\bfseries\color{Red}][2]{Modèle analytique}
    
    \point[\bfseries][6][Red]{Analyse des caches} en fonction des structures gros grain des matrices

    \subpt[][10][Red]{Analyse détaillée des \textbf{matrices bandées}}

    \smallskip

    \arrow[\bfseries\color{Red}][2]{Analyse des performances} en fonction des matrices

    \point[\bfseries][6][Red]{Critère de classification} des matrices
    
\end{frame}

\begin{frame}
    \frametitle{Perspectives}

    \textbf{\textcolor{Mulberry}{Perspectives}} Travaux à faire

    \smallskip

    \point[][2][Red]{De nombreuses \underbar{optimisations} à faire sur l'implémentation RTL}

    - Calcul des régions en matériel pour alléger le coût processeur

    - Support du calcul flottant

    - Cohérence de cache complète
    
    \smallskip

    \point[][2][Red]{Étude de l'utilisation avec d'\underbar{autres \textit{kernels}} (ex. : SpMSpM)}

    \smallskip

    Niveau logiciel

    - Librairie de kernels optimisés pour SpDCache (extension de SparseBLAS [])

    \smallskip

    \point[][2][Red]{Mono-c\oe{}ur, un niveau de cache} \arrow[]{Question du passage à l'échelle}

    \arrow[][6][Red]{\underbar{Multi-c\oe{}ur, multi-niveau} :} conceptualisé

    \smallskip

    Étudier SpDCache dans un environnement haute performance

\end{frame}

\begin{frame}[plain]

    Valentin Isaac-{-}Chassande

    valentin.isaacchassande@univ-grenoble-alpes.fr

    QRcode linkedin ?

    \bigskip

    Publications

    \point[][2]{ACM TACO}

    \point[][2]{ASAP}

    Brevet

    \point[][2]{...}

    Présentations et conférences

    \point[][2]{ASAP 2024}

    \point[][2]{HiPEAC 2026}

    \point[][2]{COMPAS ..}

    SpDCache en libre accès

    \point[][2]{...}
    
\end{frame}